\paragraph{window-\texorpdfstring{$6$}{6} 规范站点抬升的条件均匀化、局域熵势与 GFF 标度不变性}
在 window-\(6\) 的可见/隐藏 Hilbert 分解、bin-fold 规范群的直积分解、纤维多重度的三档阈值分类、\(B_3/C_3\) 根字典以及隐藏部门的 \(2,3,4\) 三质量量子化都已建立之后，还可把同一窗口上的 canonical sitewise lift 再压缩为一组完全显式的有限体积恒等式与连续极限后果。以下固定
\[
\pi:=\Fold_6^{\mathrm{bin}}:\Omega_6\longrightarrow X_6,
\qquad
F_w:=\pi^{-1}(w),
\qquad
d(w):=|F_w|.
\]
由定理 \ref{thm:xi-window6-cyclic-kernel-bc3-root-datum}、推论 \ref{cor:fold-bin-m6-fiber-hist}、定理 \ref{thm:xi-foldbin-lastbit-tail-spectrum-single-jump-staircase} 与定理 \ref{thm:xi-foldbin-gauge-representation-distribution-law-m6-spectrum} 可知：
\[
|X_6|=21,
\qquad
X_6=X_6^{\mathrm{cyc}}\sqcup X_6^{\mathrm{bdry}},
\qquad
|X_6^{\mathrm{cyc}}|=18,
\qquad
|X_6^{\mathrm{bdry}}|=3,
\]
\[
d(w)\in\{2,3,4\},
\qquad
\#\{d=2\}=8,\quad
\#\{d=3\}=4,\quad
\#\{d=4\}=9,
\]
\[
d(w)=2\iff w_6=1,
\qquad
d(w)=4\iff w_6=0,\ V_6(w)\le 8,
\qquad
d(w)=3\iff w_6=0,\ V_6(w)>8,
\]
并且
\[
G_6^{\mathrm{bin}}
\cong
\prod_{w\in X_6}\mathfrak S_{d(w)}.
\]
从现在起，对任意有限位点集 \(\Lambda\)，任取可见可容集
\[
\mathcal W_\Lambda\subset X_6^\Lambda,
\]
并定义其 canonical sitewise lift
\[
\widehat{\mathcal W}_\Lambda
:=
(\pi^\Lambda)^{-1}(\mathcal W_\Lambda)
\subset \Omega_6^\Lambda,
\qquad
\mathcal G_\Lambda:=
\bigl(G_6^{\mathrm{bin}}\bigr)^\Lambda.
\]

\begin{theorem}[规范抬升的条件均匀化]\label{thm:xi-time-part9z-window6-canonical-lift-conditional-uniformization}
设 \(\widehat\mu_\Lambda\) 是 \(\widehat{\mathcal W}_\Lambda\) 上的概率测度，并满足对一切
\[
g\in \mathcal G_\Lambda,\qquad \omega\in \widehat{\mathcal W}_\Lambda,
\]
都有
\[
\widehat\mu_\Lambda(g\omega)=\widehat\mu_\Lambda(\omega).
\]
记
\[
\mu_\Lambda:=(\pi^\Lambda)_*\widehat\mu_\Lambda.
\]
则对每个 \(w=(w_x)_{x\in\Lambda}\in \mathcal W_\Lambda\)，都有精确条件分解
\[
\widehat\mu_\Lambda(\omega\mid \pi^\Lambda=w)
=
\prod_{x\in\Lambda}
\frac{\mathbf 1_{F_{w_x}}(\omega_x)}{d(w_x)}.
\]
特别地，给定可见构型以后，全部 hidden 变量在各站点上条件独立，并且在各自纤维内严格均匀。
\end{theorem}
\begin{proof}
固定 \(w\in \mathcal W_\Lambda\)。其整条抬升纤维为
\[
E_w:=(\pi^\Lambda)^{-1}(w)=\prod_{x\in\Lambda}F_{w_x}.
\]
考虑 \(\mathcal G_\Lambda\) 中保持 \(E_w\) 的稳定子
\[
(\mathcal G_\Lambda)_w.
\]
由于 \(\mathcal G_\Lambda\) 按位点作 sitewise 作用，而
\[
G_6^{\mathrm{bin}}\cong \prod_{u\in X_6}\mathfrak S_{d(u)},
\]
故限制到第 \(x\) 个位点后，稳定子在该位点上恰给出
\[
\mathfrak S(F_{w_x})\cong \mathfrak S_{d(w_x)}.
\]
于是
\[
(\mathcal G_\Lambda)_w\cong \prod_{x\in\Lambda}\mathfrak S(F_{w_x}),
\]
它在笛卡尔积 \(E_w=\prod_x F_{w_x}\) 上显然传递：任取 \(\omega,\omega'\in E_w\)，对每个 \(x\) 可选择
\[
\sigma_x\in \mathfrak S(F_{w_x})
\]
使得 \(\sigma_x(\omega_x)=\omega'_x\)，再把这些 \(\sigma_x\) 组装即可。

由 \(\widehat\mu_\Lambda\) 对整个 \(\mathcal G_\Lambda\) 的不变性可知，其条件分布
\[
\widehat\mu_\Lambda(\,\cdot\mid \pi^\Lambda=w)
\]
对稳定子 \((\mathcal G_\Lambda)_w\) 同样不变。由于该稳定子在有限集合 \(E_w\) 上传递，故条件分布只能是 \(E_w\) 上的常数测度，即
\[
\widehat\mu_\Lambda(\omega\mid \pi^\Lambda=w)
=
\frac{\mathbf 1_{E_w}(\omega)}{|E_w|}.
\]
再由
\[
|E_w|=\prod_{x\in\Lambda}d(w_x),
\qquad
\mathbf 1_{E_w}(\omega)=\prod_{x\in\Lambda}\mathbf 1_{F_{w_x}}(\omega_x),
\]
便得到所示公式。乘积形式立即给出条件独立与各纤维内严格均匀性。证毕。
\leanverified{paper\_xi\_time\_part9z\_window6\_canonical\_lift\_conditional\_uniformization}
\end{proof}

\begin{theorem}[三层局域熵势与源项整化]\label{thm:xi-time-part9z-window6-local-entropy-source-renormalization}
设
\[
H_\Lambda:\mathcal W_\Lambda\to \RR
\]
为任意可见 Hamiltonian，并在 \(\widehat{\mathcal W}_\Lambda\) 上定义 canonical lifted Gibbs 权
\[
\widehat\mu_\Lambda(\omega)
\propto
\exp\!\bigl(-H_\Lambda(\pi^\Lambda\omega)\bigr).
\]
则 lifted 配分函数满足精确消隐公式
\[
\widehat Z_\Lambda
=
\sum_{w\in\mathcal W_\Lambda}
\exp\!\Bigl(
-H_\Lambda(w)+\sum_{x\in\Lambda}s(w_x)
\Bigr),
\qquad
s(w):=\log d(w).
\]
因此，hidden sector 对可见理论的全部新增作用恰为纯一体项
\[
H_\Lambda^{\mathrm{eff}}(w)
=
H_\Lambda(w)-\sum_{x\in\Lambda}s(w_x).
\]
在 window-\(6\) 上，
\[
s(w)=
\begin{cases}
\log 2,& w_6=1,\\[4pt]
\log 4,& w_6=0,\ V_6(w)\le 8,\\[4pt]
\log 3,& w_6=0,\ V_6(w)>8.
\end{cases}
\]

更一般地，对任意局部 observable \(h:\Omega_6\to\RR\) 与任意实源
\[
t=(t_x)_{x\in\Lambda}\in\RR^\Lambda,
\]
定义
\[
\widehat Z_\Lambda(t)
:=
\sum_{\omega\in\widehat{\mathcal W}_\Lambda}
\exp\!\Bigl(
-H_\Lambda(\pi^\Lambda\omega)+\sum_{x\in\Lambda}t_x h(\omega_x)
\Bigr).
\]
再定义
\[
M_h(t;w):=
\frac1{d(w)}\sum_{n\in F_w}e^{t h(n)}.
\]
则 \(M_h(\,\cdot\,;w)\) 对 \(t\) 为整函数，并且对实 \(t\) 可写成
\[
K_h(t;w):=\log M_h(t;w),
\]
其中 \(K_h(\,\cdot\,;w)\) 在实轴上实解析，并在 \(t=0\) 的某复邻域内可取解析对数支。此时有
\[
\widehat Z_\Lambda(t)
=
\sum_{w\in\mathcal W_\Lambda}
\exp\!\Bigl(
-H_\Lambda(w)+\sum_{x\in\Lambda}\bigl[s(w_x)+K_h(t_x;w_x)\bigr]
\Bigr).
\]
因此，消去 hidden 变量以后，不会新增任何二体、长程或非局域项；hidden sector 只留下显式的一体熵势与一体解析源项。
\end{theorem}
\begin{proof}
先看无源情形。按可见构型分组求和可得
\[
\widehat Z_\Lambda
=
\sum_{\omega\in\widehat{\mathcal W}_\Lambda}
e^{-H_\Lambda(\pi^\Lambda\omega)}
=
\sum_{w\in\mathcal W_\Lambda}
e^{-H_\Lambda(w)}\,
\bigl|(\pi^\Lambda)^{-1}(w)\bigr|.
\]
而
\[
\bigl|(\pi^\Lambda)^{-1}(w)\bigr|
=
\prod_{x\in\Lambda}|F_{w_x}|
=
\prod_{x\in\Lambda}d(w_x)
=
\exp\!\Bigl(\sum_{x\in\Lambda}s(w_x)\Bigr),
\]
故
\[
\widehat Z_\Lambda
=
\sum_{w\in\mathcal W_\Lambda}
\exp\!\Bigl(
-H_\Lambda(w)+\sum_{x\in\Lambda}s(w_x)
\Bigr).
\]
于是 hidden sector 对有效 Hamiltonian 的唯一新增贡献便是
\[
-\sum_{x\in\Lambda}s(w_x).
\]

在 window-\(6\) 上，\(s(w)\) 的三层公式直接来自
\[
d(w)=2\iff w_6=1,
\qquad
d(w)=4\iff w_6=0,\ V_6(w)\le 8,
\qquad
d(w)=3\iff w_6=0,\ V_6(w)>8.
\]

再看有源情形。同样按 \(w=\pi^\Lambda(\omega)\) 分组，
\[
\widehat Z_\Lambda(t)
=
\sum_{w\in\mathcal W_\Lambda}
e^{-H_\Lambda(w)}
\sum_{\omega:\,\pi^\Lambda\omega=w}
\exp\!\Bigl(\sum_{x\in\Lambda}t_x h(\omega_x)\Bigr).
\]
由于纤维分解为直积，
\[
(\pi^\Lambda)^{-1}(w)=\prod_{x\in\Lambda}F_{w_x},
\]
内层和相应分裂为
\[
\sum_{\omega:\,\pi^\Lambda\omega=w}
e^{\sum_x t_x h(\omega_x)}
=
\prod_{x\in\Lambda}\sum_{n\in F_{w_x}}e^{t_x h(n)}.
\]
写成
\[
\sum_{n\in F_{w_x}}e^{t_x h(n)}
=
d(w_x)\,M_h(t_x;w_x),
\]
便得到
\[
\widehat Z_\Lambda(t)
=
\sum_{w\in\mathcal W_\Lambda}
\exp\!\Bigl(
-H_\Lambda(w)+\sum_{x\in\Lambda}\bigl[s(w_x)+K_h(t_x;w_x)\bigr]
\Bigr).
\]
最后，\(M_h(\,\cdot\,;w)\) 是有限个指数函数之和，故为整函数；对实 \(t\)，其值严格为正，因此 \(K_h(\,\cdot\,;w)\) 在实轴上实解析，并在 \(t=0\) 的某复邻域内可取解析对数支。证毕。
\end{proof}

\leanverified{paper\_xi\_time\_part9z\_window6\_local\_entropy\_source\_renormalization}
\leanverified{paper\_xi\_time\_part9z\_window6\_isolated\_site\_annihilation}
\begin{theorem}[孤立站点湮没律]\label{thm:xi-time-part9z-window6-isolated-site-annihilation}
对任意 \(f:\Omega_6\to\RR\)，定义其纤维平均
\[
\bar f(w):=
\frac1{d(w)}\sum_{n\in F_w}f(n)
\]
以及纤维中心化部分
\[
f^\circ(n):=f(n)-\bar f(\pi(n)).
\]
则在任意 \(\mathcal G_\Lambda\)-不变 lift \(\widehat\mu_\Lambda\) 下，若
\[
A:\mathcal W_\Lambda\to \RR
\]
为任意有界可见 observable，并且在位点列
\[
x_1,\dots,x_m\in \Lambda
\]
中至少存在一个位点只出现一次，则
\[
\mathbf E_{\widehat\mu_\Lambda}\!\left[
A(\pi^\Lambda)\prod_{j=1}^m f^\circ(\omega_{x_j})
\right]=0.
\]
\end{theorem}
\begin{proof}
取一处在 \((x_1,\dots,x_m)\) 中恰出现一次的位点 \(x_\ast\)。记
\[
\mathcal F:=
\sigma\!\bigl(\pi^\Lambda,(\omega_y)_{y\neq x_\ast}\bigr).
\]
由定理 \ref{thm:xi-time-part9z-window6-canonical-lift-conditional-uniformization}，给定 \(\pi^\Lambda\) 以后，各站点 hidden 变量条件独立，且 \(\omega_{x_\ast}\) 在
\[
F_{(\pi^\Lambda\omega)_{x_\ast}}
\]
上严格均匀。因此
\[
\mathbf E_{\widehat\mu_\Lambda}\!\bigl[f^\circ(\omega_{x_\ast})\mid \mathcal F\bigr]
=
\frac1{d(w_{x_\ast})}\sum_{n\in F_{w_{x_\ast}}}
\bigl(f(n)-\bar f(w_{x_\ast})\bigr)
=0,
\]
其中
\[
w:=\pi^\Lambda(\omega).
\]
另一方面，因 \(x_\ast\) 仅出现一次，
\[
A(\pi^\Lambda)\prod_{j\ne j_\ast}f^\circ(\omega_{x_j})
\]
是 \(\mathcal F\)-可测。故由条件期望塔性质，
\[
\mathbf E_{\widehat\mu_\Lambda}\!\left[
A(\pi^\Lambda)\prod_{j=1}^m f^\circ(\omega_{x_j})
\right]
=0.
\]
证毕。
\end{proof}

\begin{theorem}[异点 \texorpdfstring{$n$}{n}-点函数的完全可见化]\label{thm:xi-time-part9z-window6-distinct-point-complete-visibility}
\leanverified{paper\_xi\_time\_part9z\_window6\_distinct\_point\_complete\_visibility}
设
\[
x_1,\dots,x_n\in \Lambda
\]
两两不同，且
\[
F_1,\dots,F_n:\Omega_6\to\RR
\]
为任意局部 observable。记
\[
\bar F_j(w):=\frac1{d(w)}\sum_{m\in F_w}F_j(m).
\]
则在任意 \(\mathcal G_\Lambda\)-不变 lift 下，都有
\[
\mathbf E_{\widehat\mu_\Lambda}\!\left[\prod_{j=1}^n F_j(\omega_{x_j})\right]
=
\mathbf E_{\mu_\Lambda}\!\left[\prod_{j=1}^n \bar F_j(w_{x_j})\right].
\]
特别地，一切异点截断相关函数都已完全包含在可见 sector 之中；hidden sector 对任何真正的长程 \(n\)-点结构都没有独立贡献。
\end{theorem}
\begin{proof}
把每个 \(F_j\) 分解为
\[
F_j=\bar F_j\circ \pi + F_j^\circ,
\qquad
F_j^\circ:=F_j-\bar F_j\circ\pi.
\]
将乘积
\[
\prod_{j=1}^n F_j(\omega_{x_j})
\]
展开为 \(2^n\) 项。若某一项含有至少一个中心化因子 \(F_j^\circ(\omega_{x_j})\)，由于 \(x_1,\dots,x_n\) 两两不同，该项中必存在某个位点只出现一次，于是由定理 \ref{thm:xi-time-part9z-window6-isolated-site-annihilation}，其期望为零。

因此唯一幸存项只能是纯可见项
\[
\prod_{j=1}^n \bar F_j(\pi(\omega_{x_j}))
=
\prod_{j=1}^n \bar F_j(w_{x_j}),
\]
对 \(\widehat\mu_\Lambda\) 取期望便得到
\[
\mathbf E_{\widehat\mu_\Lambda}\!\left[\prod_{j=1}^n F_j(\omega_{x_j})\right]
=
\mathbf E_{\mu_\Lambda}\!\left[\prod_{j=1}^n \bar F_j(w_{x_j})\right].
\]
证毕。
\end{proof}

\noindent
尤其，对由表 \ref{tab:fold6_b3c3_root_dictionary_b3} 与表 \ref{tab:fold6_b3c3_root_dictionary_c3} 所诱导并经 \(\pi\) 拉回到 \(\Omega_6\) 的任意 root-valued local observable，任意两两不同位点上的 \(n\)-点函数在抬升前后完全一致；因此，\(B_3/C_3\) 的长程相关结构已经在 visible quotient 中封闭。

\begin{theorem}[隐藏部门的严格局域性]\label{thm:xi-time-part9z-window6-hidden-sector-strict-locality}
\leanverified{paper\_xi\_time\_part9z\_window6\_hidden\_sector\_strict\_locality}
设
\[
h^\circ:\Omega_6\to\RR
\]
有界，并且在每条纤维上均值为零，即
\[
\frac1{d(w)}\sum_{n\in F_w}h^\circ(n)=0
\qquad
(w\in X_6).
\]
则在任意 \(\mathcal G_\Lambda\)-不变 lift 下成立：
\begin{enumerate}
  \item 对任意 \(x\neq y\) 都有
  \[
  \operatorname{Cov}_{\widehat\mu_\Lambda}\!\bigl(h^\circ(\omega_x),h^\circ(\omega_y)\bigr)=0.
  \]
  \item 令 \(\Lambda_\delta\subset D\cap \delta\ZZ^2\) 逼近有界区域 \(D\)，并定义面积型测试函数配对
  \[
  \Phi_\delta^\circ(\varphi)
  :=
  \delta^2\sum_{x\in\Lambda_\delta}\varphi(x)\,h^\circ(\omega_x),
  \qquad
  \varphi\in C_c^\infty(D).
  \]
  则
  \[
  \Var\!\bigl(\Phi_\delta^\circ(\varphi)\bigr)
  \le
  \|\varphi\|_\infty^2\|h^\circ\|_\infty^2
  \bigl(|D|\,\delta^2+o(\delta^2)\bigr),
  \]
  从而
  \[
  \Phi_\delta^\circ(\varphi)\longrightarrow 0
  \qquad\text{于 }L^2.
  \]
  \item 若改用白噪声型归一化
  \[
  \Psi_\delta^\circ(\varphi)
  :=
  \delta\sum_{x\in\Lambda_\delta}\varphi(x)\,h^\circ(\omega_x),
  \]
  则对支撑互不相交的 \(\varphi,\psi\in C_c^\infty(D)\)，有
  \[
  \operatorname{Cov}\!\bigl(\Psi_\delta^\circ(\varphi),\Psi_\delta^\circ(\psi)\bigr)=0
  \]
  对充分小的 \(\delta\) 恒成立。因而任何存在的子序列极限，其协方差泛函都只能支撑在对角线上；隐藏部门至多产生接触项或局域白噪声，而绝不可能产生 Green 核型的长程场。
\end{enumerate}
\end{theorem}
\begin{proof}
先证第 \(1\) 条。由中心化条件与定理
\ref{thm:xi-time-part9z-window6-canonical-lift-conditional-uniformization}，
\[
\mathbf E_{\widehat\mu_\Lambda}\!\bigl[h^\circ(\omega_x)\mid \pi^\Lambda\bigr]=0
\]
对每个 \(x\) 都成立，因此
\[
\mathbf E_{\widehat\mu_\Lambda}\!\bigl[h^\circ(\omega_x)\bigr]=0.
\]
当 \(x\neq y\) 时，定理
\ref{thm:xi-time-part9z-window6-isolated-site-annihilation} 在 \(m=2\)、\(A\equiv 1\)、\(f=h^\circ\) 的情形立即给出
\[
\mathbf E_{\widehat\mu_\Lambda}\!\bigl[h^\circ(\omega_x)h^\circ(\omega_y)\bigr]=0.
\]
故
\[
\operatorname{Cov}_{\widehat\mu_\Lambda}\!\bigl(h^\circ(\omega_x),h^\circ(\omega_y)\bigr)=0.
\]

再证第 \(2\) 条。由第 \(1\) 条，面积型配对的方差化为纯对角和：
\[
\Var\!\bigl(\Phi_\delta^\circ(\varphi)\bigr)
=
\delta^4\sum_{x\in\Lambda_\delta}\varphi(x)^2
\Var\!\bigl(h^\circ(\omega_x)\bigr).
\]
由于
\[
\Var\!\bigl(h^\circ(\omega_x)\bigr)\le \|h^\circ\|_\infty^2,
\qquad
\varphi(x)^2\le \|\varphi\|_\infty^2,
\]
并且
\[
|\Lambda_\delta|=|D|\,\delta^{-2}+o(\delta^{-2}),
\]
便得到
\[
\Var\!\bigl(\Phi_\delta^\circ(\varphi)\bigr)
\le
\delta^4|\Lambda_\delta|\|\varphi\|_\infty^2\|h^\circ\|_\infty^2
\le
\|\varphi\|_\infty^2\|h^\circ\|_\infty^2
\bigl(|D|\,\delta^2+o(\delta^2)\bigr).
\]
于是 \(\Phi_\delta^\circ(\varphi)\to 0\) 于 \(L^2\)。

最后看第 \(3\) 条。若 \(\operatorname{supp}\varphi\cap \operatorname{supp}\psi=\varnothing\)，则对充分小的 \(\delta\)，任何 \(x\in\Lambda_\delta\) 都不会同时落在两者支撑中。因此
\[
\operatorname{Cov}\!\bigl(\Psi_\delta^\circ(\varphi),\Psi_\delta^\circ(\psi)\bigr)
=
\delta^2\sum_{x,y\in\Lambda_\delta}
\varphi(x)\psi(y)\,
\operatorname{Cov}\!\bigl(h^\circ(\omega_x),h^\circ(\omega_y)\bigr)
\]
中，所有 \(x\neq y\) 项由第 \(1\) 条消失，而 \(x=y\) 的对角项又因支撑不交而为零，故总协方差恒为零。由此，任何可能的子序列极限都只能带有对角支撑协方差，也就是接触项或局域白噪声型行为。证毕。
\end{proof}

\endinput
